
\documentclass{article}
\usepackage{colortbl}
\usepackage{makecell}
\usepackage{multirow}
\usepackage{supertabular}

\begin{document}

\newcounter{utterance}

\centering \large Interaction Transcript for game `wordle', experiment `medium\_frequency\_words\_no\_clue\_no\_critic', episode 10 with reflexion{-}llama{-}3.1{-}8b{-}instruct.
\vspace{24pt}

{ \footnotesize  \setcounter{utterance}{1}
\setlength{\tabcolsep}{0pt}
\begin{supertabular}{c@{$\;$}|p{.15\linewidth}@{}p{.15\linewidth}p{.15\linewidth}p{.15\linewidth}p{.15\linewidth}p{.15\linewidth}}
    \# & \multicolumn{2}{c}{Player} && \multicolumn{2}{c}{Game Master} \\
    \hline

    \theutterance \stepcounter{utterance}  
    & & & \multicolumn{4}{p{0.6\linewidth}}{
        \cellcolor[rgb]{0.9,0.9,0.9}{
            \makecell[{{p{\linewidth}}}]{
                \texttt{\tiny{[P1$\langle$GM]}}
                \texttt{You have attempted this same task before and did not succeed. These are the plans you wrote after those attempts.} \\
\texttt{Trial \#1: New plan: } \\
\\ 
\texttt{I got stuck in a loop by repeatedly guessing the same word "theme" without adding any new information. This led to the same feedback for multiple guesses, and I failed to make progress towards the target word. Next time, I will avoid guessing the same word multiple times and instead try to find a new word that includes the correct letters 't', 'h', and 'e' in different positions, and excludes the letters 'm' and the duplicate 'e'. I will also pay closer attention to the feedback and adjust my strategy accordingly. For example, I will try to use the correct letters in different positions, and avoid guessing words that include the same letters multiple times.} \\
\\ 
\texttt{Trial \#2: New plan: } \\
\\ 
\texttt{I got stuck in a loop by repeatedly guessing the same word "theme" without adding any new information. This led to the same feedback for multiple guesses, and I failed to make progress towards the target word. However, I also failed to take advantage of the fact that 'h' is yellow, which means it is in the word but not in the correct position. Next time, I will try to use this information more effectively. I will try a word where 'h' is not next to 't', and I will also try to use the correct letters 't', 'h', and 'e' in different positions, and excludes the letters 'm', 'i', 'c', and the duplicate 'e'.} \\
\\ 
\texttt{Trial \#3: New plan: } \\
\\ 
\texttt{I got stuck in a loop by repeatedly guessing the same word "theme" without adding any new information, and I failed to make progress towards the target word. However, I also failed to take advantage of the fact that 'h' is yellow, which means it is in the word but not in the correct position. Furthermore, I failed to exclude the letters 's' and the duplicate 'e' after the last guess. Next time, I will try to use this information more effectively. I will try a word where 'h' is not next to 't', and I will also try to use the correct letters 't', 'h', and 'e' in different positions, and exclude the letters 'm', 'i', 'c', 's', and the duplicate 'e'. I will also pay closer attention to the feedback and adjust my strategy accordingly.} \\
\\ 
\\ 
\texttt{You are a language wizard who likes to guess words by using the given rules.} \\
\\ 
\texttt{Welcome to Wordle! You have six attempts to guess the target word, a valid English word of five lowercase letters (a{-}z). Please use the tags "explanation:" and "guess:" to provide a concise explanation for each guess.} \\
\\ 
\texttt{For instance, if your guess is "apple", your response should be} \\
\texttt{explanation: this is a common five{-}letter English word, and I am starting my guess with this word.} \\
\texttt{guess: apple} \\
\\ 
\texttt{After each guess, your answer will be validated, and you will receive feedback indicating which letters are correct (green), which letters are correct but in the wrong position (yellow), and which letters are incorrect (red). This feedback can be useful in determining which letters to include or exclude in your next guess.} \\
\\ 
\texttt{For example, the feedback for "apple" might be:} \\
\texttt{guess\_feedback: a<yellow> p<yellow> p<green> l<yellow> e<red>} \\
\\ 
\texttt{The explanation should contain details about how the guess\_feedback is used to arrive at a new guess.} \\
\\ 
\texttt{Do not generate any other text.} \\
\\ 
\texttt{Let's begin with your first guess.} \\
            }
        }
    }
     \\ \\

    \theutterance \stepcounter{utterance}  
    & \multicolumn{4}{p{0.6\linewidth}}{
        \cellcolor[rgb]{0.561,0.737,0.561}{
            \makecell[{{p{\linewidth}}}]{
                \texttt{\tiny{[P1$\rangle$GM]}}
                \texttt{explanation: I will start with a common five{-}letter English word that includes the letters 't', 'h', and 'e' in different positions, and excludes the letters 'm' and the duplicate 'e'. } \\
\texttt{guess: theme} \\
            }
        }
    }
    & & \\ \\

    \theutterance \stepcounter{utterance}  
    & & & \multicolumn{2}{p{0.3\linewidth}}{
        \cellcolor[rgb]{0.9,0.9,0.9}{
            \makecell[{{p{\linewidth}}}]{
                \texttt{\tiny{[GM$|$GM]}}
                \texttt{attempts = 1} \\
\texttt{target\_word = tenth} \\
\texttt{guess = theme} \\
\texttt{guess\_feedback = t<green> h<yellow> e<yellow> m<red> e<red>} \\
            }
        }
    }
    & & \\ \\

    \theutterance \stepcounter{utterance}  
    & & & \multicolumn{4}{p{0.6\linewidth}}{
        \cellcolor[rgb]{0.9,0.9,0.9}{
            \makecell[{{p{\linewidth}}}]{
                \texttt{\tiny{[P1$\langle$GM]}}
                \texttt{guess\_feedback: t<green> h<yellow> e<yellow> m<red> e<red>} \\
\\ 
\texttt{Provide your response only in this format.} \\
\texttt{explanation: details} \\
\texttt{guess: word} \\
            }
        }
    }
     \\ \\

    \theutterance \stepcounter{utterance}  
    & \multicolumn{4}{p{0.6\linewidth}}{
        \cellcolor[rgb]{0.561,0.737,0.561}{
            \makecell[{{p{\linewidth}}}]{
                \texttt{\tiny{[P1$\rangle$GM]}}
                \texttt{explanation: Since 't' is green, it is in the correct position. 'h' and 'e' are yellow, meaning they are in the word but not in the correct position. I will try to use 'h' and 'e' in different positions and exclude the letters 'm' and the duplicate 'e'. } \\
\texttt{guess: theme} \\
            }
        }
    }
    & & \\ \\

    \theutterance \stepcounter{utterance}  
    & & & \multicolumn{2}{p{0.3\linewidth}}{
        \cellcolor[rgb]{0.9,0.9,0.9}{
            \makecell[{{p{\linewidth}}}]{
                \texttt{\tiny{[GM$|$GM]}}
                \texttt{attempts = 2} \\
\texttt{target\_word = tenth} \\
\texttt{guess = theme} \\
\texttt{guess\_feedback = t<green> h<yellow> e<yellow> m<red> e<red>} \\
            }
        }
    }
    & & \\ \\

    \theutterance \stepcounter{utterance}  
    & & & \multicolumn{4}{p{0.6\linewidth}}{
        \cellcolor[rgb]{0.9,0.9,0.9}{
            \makecell[{{p{\linewidth}}}]{
                \texttt{\tiny{[P1$\langle$GM]}}
                \texttt{guess\_feedback: t<green> h<yellow> e<yellow> m<red> e<red>} \\
\\ 
\texttt{Provide your response only in this format.} \\
\texttt{explanation: details} \\
\texttt{guess: word} \\
            }
        }
    }
     \\ \\

    \theutterance \stepcounter{utterance}  
    & \multicolumn{4}{p{0.6\linewidth}}{
        \cellcolor[rgb]{0.561,0.737,0.561}{
            \makecell[{{p{\linewidth}}}]{
                \texttt{\tiny{[P1$\rangle$GM]}}
                \texttt{explanation: I noticed that the feedback for the second guess is the same as the first guess, indicating that I'm stuck in a loop. To avoid this, I will try a new word that includes the correct letters 't', 'h', and 'e' in different positions, and excludes the letters 'm' and the duplicate 'e'. I will also try to use 'h' in a different position than 't'.} \\
\texttt{guess: theme} \\
            }
        }
    }
    & & \\ \\

    \theutterance \stepcounter{utterance}  
    & & & \multicolumn{2}{p{0.3\linewidth}}{
        \cellcolor[rgb]{0.9,0.9,0.9}{
            \makecell[{{p{\linewidth}}}]{
                \texttt{\tiny{[GM$|$GM]}}
                \texttt{attempts = 3} \\
\texttt{target\_word = tenth} \\
\texttt{guess = theme} \\
\texttt{guess\_feedback = t<green> h<yellow> e<yellow> m<red> e<red>} \\
            }
        }
    }
    & & \\ \\

    \theutterance \stepcounter{utterance}  
    & & & \multicolumn{4}{p{0.6\linewidth}}{
        \cellcolor[rgb]{0.9,0.9,0.9}{
            \makecell[{{p{\linewidth}}}]{
                \texttt{\tiny{[P1$\langle$GM]}}
                \texttt{guess\_feedback: t<green> h<yellow> e<yellow> m<red> e<red>} \\
\\ 
\texttt{Provide your response only in this format.} \\
\texttt{explanation: details} \\
\texttt{guess: word} \\
            }
        }
    }
     \\ \\

    \theutterance \stepcounter{utterance}  
    & \multicolumn{4}{p{0.6\linewidth}}{
        \cellcolor[rgb]{0.561,0.737,0.561}{
            \makecell[{{p{\linewidth}}}]{
                \texttt{\tiny{[P1$\rangle$GM]}}
                \texttt{explanation: Since the feedback is still the same, I will try a new approach. I will keep the correct letter 't' in the correct position and try to use 'h' and 'e' in different positions. I will also exclude the letters 'm' and the duplicate 'e'. } \\
\texttt{guess: theft} \\
            }
        }
    }
    & & \\ \\

    \theutterance \stepcounter{utterance}  
    & & & \multicolumn{2}{p{0.3\linewidth}}{
        \cellcolor[rgb]{0.9,0.9,0.9}{
            \makecell[{{p{\linewidth}}}]{
                \texttt{\tiny{[GM$|$GM]}}
                \texttt{attempts = 4} \\
\texttt{target\_word = tenth} \\
\texttt{guess = theft} \\
\texttt{guess\_feedback = t<green> h<yellow> e<yellow> f<red> t<yellow>} \\
            }
        }
    }
    & & \\ \\

    \theutterance \stepcounter{utterance}  
    & & & \multicolumn{4}{p{0.6\linewidth}}{
        \cellcolor[rgb]{0.9,0.9,0.9}{
            \makecell[{{p{\linewidth}}}]{
                \texttt{\tiny{[P1$\langle$GM]}}
                \texttt{guess\_feedback: t<green> h<yellow> e<yellow> f<red> t<yellow>} \\
\\ 
\texttt{Provide your response only in this format.} \\
\texttt{explanation: details} \\
\texttt{guess: word} \\
            }
        }
    }
     \\ \\

    \theutterance \stepcounter{utterance}  
    & \multicolumn{4}{p{0.6\linewidth}}{
        \cellcolor[rgb]{0.561,0.737,0.561}{
            \makecell[{{p{\linewidth}}}]{
                \texttt{\tiny{[P1$\rangle$GM]}}
                \texttt{explanation: The feedback shows that 't' is still in the correct position, and 'h' and 'e' are in the word but not in the correct position. 'f' is incorrect, and 't' is also yellow, indicating that it's in the word but not in the correct position. I will try to use 'h' and 'e' in different positions and find a new position for 't'.} \\
\texttt{guess: theme} \\
            }
        }
    }
    & & \\ \\

    \theutterance \stepcounter{utterance}  
    & & & \multicolumn{2}{p{0.3\linewidth}}{
        \cellcolor[rgb]{0.9,0.9,0.9}{
            \makecell[{{p{\linewidth}}}]{
                \texttt{\tiny{[GM$|$GM]}}
                \texttt{attempts = 5} \\
\texttt{target\_word = tenth} \\
\texttt{guess = theme} \\
\texttt{guess\_feedback = t<green> h<yellow> e<yellow> m<red> e<red>} \\
            }
        }
    }
    & & \\ \\

    \theutterance \stepcounter{utterance}  
    & & & \multicolumn{4}{p{0.6\linewidth}}{
        \cellcolor[rgb]{0.9,0.9,0.9}{
            \makecell[{{p{\linewidth}}}]{
                \texttt{\tiny{[P1$\langle$GM]}}
                \texttt{guess\_feedback: t<green> h<yellow> e<yellow> m<red> e<red>} \\
\\ 
\texttt{Provide your response only in this format.} \\
\texttt{explanation: details} \\
\texttt{guess: word} \\
            }
        }
    }
     \\ \\

    \theutterance \stepcounter{utterance}  
    & \multicolumn{4}{p{0.6\linewidth}}{
        \cellcolor[rgb]{0.561,0.737,0.561}{
            \makecell[{{p{\linewidth}}}]{
                \texttt{\tiny{[P1$\rangle$GM]}}
                \texttt{explanation: The feedback is the same as before, indicating that I'm stuck in a loop. I will try a new word that includes the correct letters 't', 'h', and 'e' in different positions, and excludes the letters 'm' and the duplicate 'e'. I will also try to use the new information that 't' is in the correct position and 't' is also yellow in the previous guess.} \\
\texttt{guess: their} \\
            }
        }
    }
    & & \\ \\

    \theutterance \stepcounter{utterance}  
    & & & \multicolumn{2}{p{0.3\linewidth}}{
        \cellcolor[rgb]{0.9,0.9,0.9}{
            \makecell[{{p{\linewidth}}}]{
                \texttt{\tiny{[GM$|$GM]}}
                \texttt{attempts = 6} \\
\texttt{target\_word = tenth} \\
\texttt{guess = their} \\
\texttt{guess\_feedback = t<green> h<yellow> e<yellow> i<red> r<red>} \\
            }
        }
    }
    & & \\ \\

    \theutterance \stepcounter{utterance}  
    & & & \multicolumn{2}{p{0.3\linewidth}}{
        \cellcolor[rgb]{0.9,0.9,0.9}{
            \makecell[{{p{\linewidth}}}]{
                \texttt{\tiny{[GM$|$GM]}}
                \texttt{game\_result = LOSS} \\
            }
        }
    }
    & & \\ \\

    \theutterance \stepcounter{utterance}  
    & & \multicolumn{4}{p{0.6\linewidth}}{
        \cellcolor[rgb]{0.95,0.95,0.95}{
            \makecell[{{p{\linewidth}}}]{
                \texttt{\tiny{[P1$\leftrightarrow$P]}}
                \texttt{Post{-}episode Reflexion (not sent to GM):} \\
\texttt{New plan: } \\
\\ 
\texttt{I got stuck in a loop by repeatedly guessing the same word "theme" without adding any new information, and I failed to make progress towards the target word. However, I also failed to take advantage of the fact that 'h' is yellow, which means it is in the word but not in the correct position. Furthermore, I failed to exclude the letters 's' and the duplicate 'e' after the last guess. Next time, I will try to use this information more effectively. I will try a word where 'h' is not next to 't', and I will also try to use the correct letters 't', 'h', and 'e' in different positions, and exclude the letters 'm', 'i', 'c', 's', and the duplicate 'e'. I will also pay closer attention to the feedback and adjust my strategy accordingly. Additionally, I will try to use the fact that 't' is both green and yellow in the previous guesses to my advantage, and try to find a word where 't' is in a different position than in the previous guesses.} \\
            }
        }
    }
    & \\ \\

\end{supertabular}
}

\end{document}
